% theory_main.tex — Main paper theory subsection (~0.5 page)
% Include with % theory_main.tex — Main paper theory subsection (~0.5 page)
% Include with % theory_main.tex — Main paper theory subsection (~0.5 page)
% Include with % theory_main.tex — Main paper theory subsection (~0.5 page)
% Include with \input{theory_main} in the appropriate section of paper.tex

\subsection{Theoretical Guarantees}
\label{sec:theory}

We now formalise when JEPA pretraining retains event-relevant information.
All proofs and assumptions are in \cref{app:theory}.

\paragraph{Setup.}
Let $X_{\leq t}$ denote observations up to time $t$, and let $E_{t+\Delta t} \in \{0,1\}$ indicate whether an event occurs in the interval $(t, t+\Delta t]$.
The encoder produces $H_t = \encoder(X_{\leq t}) \in \R^d$, and the target encoder produces $H^* = \target(X_{(t,t+\Delta t]}) \in \R^d$ from the future interval.
We write $\hat{H} = \predictor(H_t, \Delta t)$ for the predicted representation.
We use $I(\cdot;\cdot)$ for mutual information and $H(\cdot)$ for entropy, following \citet{cover2006elements}.

\begin{definition}[Predictive and event-relevant information]
\label{def:info_decomposition}
The \emph{total predictive information} in $H_t$ about the future interval is $I(H_t; X_{(t,t+\Delta t]})$.
The \emph{event-relevant information} captured by the target encoder is $I(H^*; E_{t+\Delta t})$.
The \emph{event-irrelevant} component of $H^*$ is the residual $I(H^*; X_{(t,t+\Delta t]} \mid E_{t+\Delta t})$.
\end{definition}

\begin{proposition}[Event-Information Retention]
\label{prop:retention}
Suppose \textup{(A1)} the event $E_{t+\Delta t}$ is conditionally independent of $X_{\leq t}$ given $H^*$,
\textup{(A2)} the pretraining loss satisfies $\E[\|\hat{H} - H^*\|_2^2] \leq \varepsilon$,
\textup{(A3)} the event posterior $\eta(h) = P(E_{t+\Delta t}=1 \mid H^*=h)$ is $L$-Lipschitz,
and \textup{(A4)} the marginal event rate satisfies $p = P(E_{t+\Delta t}=1) \in [\underline{p},\, \overline{p}] \subset (0,1)$.
Then
\begin{equation}
\label{eq:main_bound}
I(H_t;\, E_{t+\Delta t}) \;\geq\; I(H^*;\, E_{t+\Delta t}) \;-\; C_p \, L^2 \, \varepsilon\;,
\end{equation}
where $C_p = \bigl(2\,\underline{p}\,(1-\overline{p})\bigr)^{-1}$ depends only on the event-rate bounds.
\end{proposition}

\begin{proof}[Proof sketch]
Since $\hat{H}$ is a deterministic function of $(H_t, \Delta t)$, the data processing inequality gives $I(H_t; E) \geq I(\hat{H}; E)$.
Under A1, $P(E{=}1 \mid \hat{H}) = \E[\eta(H^*) \mid \hat{H}]$.
A Jensen-gap argument on the convex KL divergence, combined with the Lipschitz bound A3 and prediction error A2, yields $I(H^*; E) - I(\hat{H}; E) \leq C_p L^2 \varepsilon$.
Full proof in \cref{app:theory}.
\end{proof}

\begin{corollary}[Precursor necessity]
\label{cor:precursor}
The bound \eqref{eq:main_bound} is non-vacuous (i.e.\ the right-hand side is positive) if and only if:
\begin{enumerate}[nosep,label=(\roman*)]
  \item $I(H^*; E_{t+\Delta t}) > 0$: the future interval contains event precursors that the target encoder captures;
  \item $\varepsilon < I(H^*; E_{t+\Delta t}) / (C_p L^2)$: the predictor approximates the target representation well enough;
  \item The encoder $\encoder$ has sufficient capacity to represent the precursors (implicit in A1).
\end{enumerate}
\end{corollary}

\paragraph{Connection to experiments.}
\Cref{cor:precursor} directly explains our empirical findings.
On C-MAPSS (\cref{sec:experiments}), degradation unfolds over hundreds of cycles, so $I(H^*; E_{t+\Delta t})$ is large and $\varepsilon$ is driven small by pretraining, yielding strong downstream performance (h-AUROC $\geq 0.74$).
On CHB-MIT, seizure onset in scalp EEG has minimal precursors in a 16-second context, so $I(H^*; E_{t+\Delta t}) \approx 0$ regardless of $\varepsilon$: condition (i) fails.
That predictor \emph{weights} do not transfer (\cref{sec:finetune_ablation}) while the encoder does is consistent with the bound: the encoder preserves information (the $I(H_t; E_{t+\Delta t})$ term), but the predictor's pretraining codomain ($\R^d$ representations) differs from the downstream codomain ($[0,1]$ event probabilities), so pretrained weights offer no inductive bias for the finetuning objective.

\paragraph{Label efficiency.}
When $I(H_t; E_{t+\Delta t})$ is large, very few labeled examples suffice to steer the predictor toward the event.
Formally (see \cref{app:theory}), fixing the encoder and optimising the head with $n$ labeled samples, the expected excess risk is bounded by a \emph{representation gap} term that depends only on $I(H_t; E_{t+\Delta t})$ plus a \emph{statistical error} term $O(1/\sqrt{n})$.
For C-MAPSS FD001, where $I(H_t; E_{t+\Delta t})$ is large, the statistical error term already falls below the representation gap by $n \approx 200$ windows.
This is consistent with our empirical observation that FAM retains $\mathbf{92\%}$ of full-label h-AUROC using just 2 training engines out of 85 (2\% labels; \cref{tab:sub5pct}).
 in the appropriate section of paper.tex

\subsection{Theoretical Guarantees}
\label{sec:theory}

We now formalise when JEPA pretraining retains event-relevant information.
All proofs and assumptions are in \cref{app:theory}.

\paragraph{Setup.}
Let $X_{\leq t}$ denote observations up to time $t$, and let $E_{t+\Delta t} \in \{0,1\}$ indicate whether an event occurs in the interval $(t, t+\Delta t]$.
The encoder produces $H_t = \encoder(X_{\leq t}) \in \R^d$, and the target encoder produces $H^* = \target(X_{(t,t+\Delta t]}) \in \R^d$ from the future interval.
We write $\hat{H} = \predictor(H_t, \Delta t)$ for the predicted representation.
We use $I(\cdot;\cdot)$ for mutual information and $H(\cdot)$ for entropy, following \citet{cover2006elements}.

\begin{definition}[Predictive and event-relevant information]
\label{def:info_decomposition}
The \emph{total predictive information} in $H_t$ about the future interval is $I(H_t; X_{(t,t+\Delta t]})$.
The \emph{event-relevant information} captured by the target encoder is $I(H^*; E_{t+\Delta t})$.
The \emph{event-irrelevant} component of $H^*$ is the residual $I(H^*; X_{(t,t+\Delta t]} \mid E_{t+\Delta t})$.
\end{definition}

\begin{proposition}[Event-Information Retention]
\label{prop:retention}
Suppose \textup{(A1)} the event $E_{t+\Delta t}$ is conditionally independent of $X_{\leq t}$ given $H^*$,
\textup{(A2)} the pretraining loss satisfies $\E[\|\hat{H} - H^*\|_2^2] \leq \varepsilon$,
\textup{(A3)} the event posterior $\eta(h) = P(E_{t+\Delta t}=1 \mid H^*=h)$ is $L$-Lipschitz,
and \textup{(A4)} the marginal event rate satisfies $p = P(E_{t+\Delta t}=1) \in [\underline{p},\, \overline{p}] \subset (0,1)$.
Then
\begin{equation}
\label{eq:main_bound}
I(H_t;\, E_{t+\Delta t}) \;\geq\; I(H^*;\, E_{t+\Delta t}) \;-\; C_p \, L^2 \, \varepsilon\;,
\end{equation}
where $C_p = \bigl(2\,\underline{p}\,(1-\overline{p})\bigr)^{-1}$ depends only on the event-rate bounds.
\end{proposition}

\begin{proof}[Proof sketch]
Since $\hat{H}$ is a deterministic function of $(H_t, \Delta t)$, the data processing inequality gives $I(H_t; E) \geq I(\hat{H}; E)$.
Under A1, $P(E{=}1 \mid \hat{H}) = \E[\eta(H^*) \mid \hat{H}]$.
A Jensen-gap argument on the convex KL divergence, combined with the Lipschitz bound A3 and prediction error A2, yields $I(H^*; E) - I(\hat{H}; E) \leq C_p L^2 \varepsilon$.
Full proof in \cref{app:theory}.
\end{proof}

\begin{corollary}[Precursor necessity]
\label{cor:precursor}
The bound \eqref{eq:main_bound} is non-vacuous (i.e.\ the right-hand side is positive) if and only if:
\begin{enumerate}[nosep,label=(\roman*)]
  \item $I(H^*; E_{t+\Delta t}) > 0$: the future interval contains event precursors that the target encoder captures;
  \item $\varepsilon < I(H^*; E_{t+\Delta t}) / (C_p L^2)$: the predictor approximates the target representation well enough;
  \item The encoder $\encoder$ has sufficient capacity to represent the precursors (implicit in A1).
\end{enumerate}
\end{corollary}

\paragraph{Connection to experiments.}
\Cref{cor:precursor} directly explains our empirical findings.
On C-MAPSS (\cref{sec:experiments}), degradation unfolds over hundreds of cycles, so $I(H^*; E_{t+\Delta t})$ is large and $\varepsilon$ is driven small by pretraining, yielding strong downstream performance (h-AUROC $\geq 0.74$).
On CHB-MIT, seizure onset in scalp EEG has minimal precursors in a 16-second context, so $I(H^*; E_{t+\Delta t}) \approx 0$ regardless of $\varepsilon$: condition (i) fails.
That predictor \emph{weights} do not transfer (\cref{sec:finetune_ablation}) while the encoder does is consistent with the bound: the encoder preserves information (the $I(H_t; E_{t+\Delta t})$ term), but the predictor's pretraining codomain ($\R^d$ representations) differs from the downstream codomain ($[0,1]$ event probabilities), so pretrained weights offer no inductive bias for the finetuning objective.

\paragraph{Label efficiency.}
When $I(H_t; E_{t+\Delta t})$ is large, very few labeled examples suffice to steer the predictor toward the event.
Formally (see \cref{app:theory}), fixing the encoder and optimising the head with $n$ labeled samples, the expected excess risk is bounded by a \emph{representation gap} term that depends only on $I(H_t; E_{t+\Delta t})$ plus a \emph{statistical error} term $O(1/\sqrt{n})$.
For C-MAPSS FD001, where $I(H_t; E_{t+\Delta t})$ is large, the statistical error term already falls below the representation gap by $n \approx 200$ windows.
This is consistent with our empirical observation that FAM retains $\mathbf{92\%}$ of full-label h-AUROC using just 2 training engines out of 85 (2\% labels; \cref{tab:sub5pct}).
 in the appropriate section of paper.tex

\subsection{Theoretical Guarantees}
\label{sec:theory}

We now formalise when JEPA pretraining retains event-relevant information.
All proofs and assumptions are in \cref{app:theory}.

\paragraph{Setup.}
Let $X_{\leq t}$ denote observations up to time $t$, and let $E_{t+\Delta t} \in \{0,1\}$ indicate whether an event occurs in the interval $(t, t+\Delta t]$.
The encoder produces $H_t = \encoder(X_{\leq t}) \in \R^d$, and the target encoder produces $H^* = \target(X_{(t,t+\Delta t]}) \in \R^d$ from the future interval.
We write $\hat{H} = \predictor(H_t, \Delta t)$ for the predicted representation.
We use $I(\cdot;\cdot)$ for mutual information and $H(\cdot)$ for entropy, following \citet{cover2006elements}.

\begin{definition}[Predictive and event-relevant information]
\label{def:info_decomposition}
The \emph{total predictive information} in $H_t$ about the future interval is $I(H_t; X_{(t,t+\Delta t]})$.
The \emph{event-relevant information} captured by the target encoder is $I(H^*; E_{t+\Delta t})$.
The \emph{event-irrelevant} component of $H^*$ is the residual $I(H^*; X_{(t,t+\Delta t]} \mid E_{t+\Delta t})$.
\end{definition}

\begin{proposition}[Event-Information Retention]
\label{prop:retention}
Suppose \textup{(A1)} the event $E_{t+\Delta t}$ is conditionally independent of $X_{\leq t}$ given $H^*$,
\textup{(A2)} the pretraining loss satisfies $\E[\|\hat{H} - H^*\|_2^2] \leq \varepsilon$,
\textup{(A3)} the event posterior $\eta(h) = P(E_{t+\Delta t}=1 \mid H^*=h)$ is $L$-Lipschitz,
and \textup{(A4)} the marginal event rate satisfies $p = P(E_{t+\Delta t}=1) \in [\underline{p},\, \overline{p}] \subset (0,1)$.
Then
\begin{equation}
\label{eq:main_bound}
I(H_t;\, E_{t+\Delta t}) \;\geq\; I(H^*;\, E_{t+\Delta t}) \;-\; C_p \, L^2 \, \varepsilon\;,
\end{equation}
where $C_p = \bigl(2\,\underline{p}\,(1-\overline{p})\bigr)^{-1}$ depends only on the event-rate bounds.
\end{proposition}

\begin{proof}[Proof sketch]
Since $\hat{H}$ is a deterministic function of $(H_t, \Delta t)$, the data processing inequality gives $I(H_t; E) \geq I(\hat{H}; E)$.
Under A1, $P(E{=}1 \mid \hat{H}) = \E[\eta(H^*) \mid \hat{H}]$.
A Jensen-gap argument on the convex KL divergence, combined with the Lipschitz bound A3 and prediction error A2, yields $I(H^*; E) - I(\hat{H}; E) \leq C_p L^2 \varepsilon$.
Full proof in \cref{app:theory}.
\end{proof}

\begin{corollary}[Precursor necessity]
\label{cor:precursor}
The bound \eqref{eq:main_bound} is non-vacuous (i.e.\ the right-hand side is positive) if and only if:
\begin{enumerate}[nosep,label=(\roman*)]
  \item $I(H^*; E_{t+\Delta t}) > 0$: the future interval contains event precursors that the target encoder captures;
  \item $\varepsilon < I(H^*; E_{t+\Delta t}) / (C_p L^2)$: the predictor approximates the target representation well enough;
  \item The encoder $\encoder$ has sufficient capacity to represent the precursors (implicit in A1).
\end{enumerate}
\end{corollary}

\paragraph{Connection to experiments.}
\Cref{cor:precursor} directly explains our empirical findings.
On C-MAPSS (\cref{sec:experiments}), degradation unfolds over hundreds of cycles, so $I(H^*; E_{t+\Delta t})$ is large and $\varepsilon$ is driven small by pretraining, yielding strong downstream performance (h-AUROC $\geq 0.74$).
On CHB-MIT, seizure onset in scalp EEG has minimal precursors in a 16-second context, so $I(H^*; E_{t+\Delta t}) \approx 0$ regardless of $\varepsilon$: condition (i) fails.
That predictor \emph{weights} do not transfer (\cref{sec:finetune_ablation}) while the encoder does is consistent with the bound: the encoder preserves information (the $I(H_t; E_{t+\Delta t})$ term), but the predictor's pretraining codomain ($\R^d$ representations) differs from the downstream codomain ($[0,1]$ event probabilities), so pretrained weights offer no inductive bias for the finetuning objective.

\paragraph{Label efficiency.}
When $I(H_t; E_{t+\Delta t})$ is large, very few labeled examples suffice to steer the predictor toward the event.
Formally (see \cref{app:theory}), fixing the encoder and optimising the head with $n$ labeled samples, the expected excess risk is bounded by a \emph{representation gap} term that depends only on $I(H_t; E_{t+\Delta t})$ plus a \emph{statistical error} term $O(1/\sqrt{n})$.
For C-MAPSS FD001, where $I(H_t; E_{t+\Delta t})$ is large, the statistical error term already falls below the representation gap by $n \approx 200$ windows.
This is consistent with our empirical observation that FAM retains $\mathbf{92\%}$ of full-label h-AUROC using just 2 training engines out of 85 (2\% labels; \cref{tab:sub5pct}).
 in the appropriate section of paper.tex

\subsection{Theoretical Guarantees}
\label{sec:theory}

We now formalise when JEPA pretraining retains event-relevant information.
All proofs and assumptions are in \cref{app:theory}.

\paragraph{Setup.}
Let $X_{\leq t}$ denote observations up to time $t$, and let $E_{t+\Delta t} \in \{0,1\}$ indicate whether an event occurs in the interval $(t, t+\Delta t]$.
The encoder produces $H_t = \encoder(X_{\leq t}) \in \R^d$, and the target encoder produces $H^* = \target(X_{(t,t+\Delta t]}) \in \R^d$ from the future interval.
We write $\hat{H} = \predictor(H_t, \Delta t)$ for the predicted representation.
We use $I(\cdot;\cdot)$ for mutual information and $H(\cdot)$ for entropy, following \citet{cover2006elements}.

\begin{definition}[Predictive and event-relevant information]
\label{def:info_decomposition}
The \emph{total predictive information} in $H_t$ about the future interval is $I(H_t; X_{(t,t+\Delta t]})$.
The \emph{event-relevant information} captured by the target encoder is $I(H^*; E_{t+\Delta t})$.
The \emph{event-irrelevant} component of $H^*$ is the residual $I(H^*; X_{(t,t+\Delta t]} \mid E_{t+\Delta t})$.
\end{definition}

\begin{proposition}[Event-Information Retention]
\label{prop:retention}
Suppose \textup{(A1)} the event $E_{t+\Delta t}$ is conditionally independent of $X_{\leq t}$ given $H^*$,
\textup{(A2)} the pretraining loss satisfies $\E[\|\hat{H} - H^*\|_2^2] \leq \varepsilon$,
\textup{(A3)} the event posterior $\eta(h) = P(E_{t+\Delta t}=1 \mid H^*=h)$ is $L$-Lipschitz,
and \textup{(A4)} the marginal event rate satisfies $p = P(E_{t+\Delta t}=1) \in [\underline{p},\, \overline{p}] \subset (0,1)$.
Then
\begin{equation}
\label{eq:main_bound}
I(H_t;\, E_{t+\Delta t}) \;\geq\; I(H^*;\, E_{t+\Delta t}) \;-\; C_p \, L^2 \, \varepsilon\;,
\end{equation}
where $C_p = \bigl(2\,\underline{p}\,(1-\overline{p})\bigr)^{-1}$ depends only on the event-rate bounds.
\end{proposition}

\begin{proof}[Proof sketch]
Since $\hat{H}$ is a deterministic function of $(H_t, \Delta t)$, the data processing inequality gives $I(H_t; E) \geq I(\hat{H}; E)$.
Under A1, $P(E{=}1 \mid \hat{H}) = \E[\eta(H^*) \mid \hat{H}]$.
A Jensen-gap argument on the convex KL divergence, combined with the Lipschitz bound A3 and prediction error A2, yields $I(H^*; E) - I(\hat{H}; E) \leq C_p L^2 \varepsilon$.
Full proof in \cref{app:theory}.
\end{proof}

\begin{corollary}[Precursor necessity]
\label{cor:precursor}
The bound \eqref{eq:main_bound} is non-vacuous (i.e.\ the right-hand side is positive) if and only if:
\begin{enumerate}[nosep,label=(\roman*)]
  \item $I(H^*; E_{t+\Delta t}) > 0$: the future interval contains event precursors that the target encoder captures;
  \item $\varepsilon < I(H^*; E_{t+\Delta t}) / (C_p L^2)$: the predictor approximates the target representation well enough;
  \item The encoder $\encoder$ has sufficient capacity to represent the precursors (implicit in A1).
\end{enumerate}
\end{corollary}

\paragraph{Connection to experiments.}
\Cref{cor:precursor} directly explains our empirical findings.
On C-MAPSS (\cref{sec:experiments}), degradation unfolds over hundreds of cycles, so $I(H^*; E_{t+\Delta t})$ is large and $\varepsilon$ is driven small by pretraining, yielding strong downstream performance (h-AUROC $\geq 0.74$).
On CHB-MIT, seizure onset in scalp EEG has minimal precursors in a 16-second context, so $I(H^*; E_{t+\Delta t}) \approx 0$ regardless of $\varepsilon$: condition (i) fails.
That predictor \emph{weights} do not transfer (\cref{sec:finetune_ablation}) while the encoder does is consistent with the bound: the encoder preserves information (the $I(H_t; E_{t+\Delta t})$ term), but the predictor's pretraining codomain ($\R^d$ representations) differs from the downstream codomain ($[0,1]$ event probabilities), so pretrained weights offer no inductive bias for the finetuning objective.

\paragraph{Label efficiency.}
When $I(H_t; E_{t+\Delta t})$ is large, very few labeled examples suffice to steer the predictor toward the event.
Formally (see \cref{app:theory}), fixing the encoder and optimising the head with $n$ labeled samples, the expected excess risk is bounded by a \emph{representation gap} term that depends only on $I(H_t; E_{t+\Delta t})$ plus a \emph{statistical error} term $O(1/\sqrt{n})$.
For C-MAPSS FD001, where $I(H_t; E_{t+\Delta t})$ is large, the statistical error term already falls below the representation gap by $n \approx 200$ windows.
This is consistent with our empirical observation that FAM retains $\mathbf{92\%}$ of full-label h-AUROC using just 2 training engines out of 85 (2\% labels; \cref{tab:sub5pct}).
